\begin{example}
\label{example-noetherian-topology-on-spec}
\begin{reference}
Comment by Lukas Heger of November 12, 2020.
\end{reference}
Let $R$ be a ring and $X = \Spec(R)$. Since closed subsets of $X$ correspond to
radical ideas of $R$ (Lemma \ref{lemma-Zariski-topology}) we see that $X$ is a
Noetherian topological space if and only if we have ACC for radical ideals.
This holds if and only if every radical ideal is the radical of a finitely
generated ideal (details omitted). Let
$$
\mathcal{F} = \{I \subset R \mid
\sqrt{I} = \sqrt{(f_1, \ldots, f_n)}\text{ for some }n
\text{ and }f_1, \ldots, f_n \in R\}.
$$
The reader can show that $\mathcal{F}$ is an Oka family by using the identity
$$
\sqrt{I} = \sqrt{(I, a)(I : a)}
$$
which holds for any ideal $I \subset R$ and any element $a \in R$.
On the other hand, if we have a totally ordered chain of ideals
$\{I_\alpha\}$ none of which are in $\mathcal{F}$, then the union
$I = \bigcup I_\alpha$ cannot be in $\mathcal{F}$ either. Otherwise
$\sqrt{I} = \sqrt{(f_1, \ldots, f_n)}$, then $f_i^e \in I$ for some $e$,
then $f_i^e \in I_\alpha$ for some $\alpha$ independent of $i$, then
$\sqrt{I_\alpha} = \sqrt{(f_1, \ldots, f_n)}$, contradiction.
Thus if the set of ideals not in $\mathcal{F}$ is nonempty, then it
has maximal elements and
exactly as in Lemma \ref{lemma-cohen} we conclude that $X$ is a
Noetherian topological space if and only if every prime ideal of $R$
is equal to $\sqrt{(f_1, \ldots, f_n)}$ for some $f_1, \ldots, f_n \in R$.
If we ever need this result we will carefully state and prove this
result here.
\end{example}